\chapter{Lezione 2 -- Modelli, astrazioni e link}

\section{Di cosa parliamo}

Un \textbf{sistema distribuito} è un insieme di computer (li chiameremo \textbf{processi}) che collaborano scambiandosi messaggi su una rete. Nessuno vede lo stato degli altri: ognuno sa solo cosa ha in memoria e quali messaggi ha ricevuto.

Il corso lavora su \textbf{modelli}: descrizioni matematiche semplificate del sistema, precise abbastanza da poter \textbf{dimostrare} che un algoritmo funziona. Questa lezione costruisce il modello di base (processi, esecuzioni, guasti), spiega come si scrivono le specifiche (safety e liveness) e chiude con i primi algoritmi: come costruire canali di comunicazione affidabili partendo da una rete inaffidabile.

\section{Il sistema e i processi}

\subsection{Il sistema}

\begin{inbreve}

Abbiamo \(n\) processi, ognuno con un nome diverso, collegati da canali su cui si mandano messaggi.

\end{inbreve}

\begin{itemize}
\tightlist
\item
  \(\Pi = \{p_0, \dots, p_{n-1}\}\): l'insieme degli \(n\) processi, con \textbf{identità distinte} (ognuno sa chi è e conosce i nomi degli altri).
\item
  I processi sono collegati da un \textbf{grafo di comunicazione} \(G\): i nodi sono i processi, gli archi sono i \textbf{link} (canali). Di solito \(G\) è \textbf{completo}: tutti possono parlare direttamente con tutti.
\end{itemize}

\subsection{Il processo}

\begin{inbreve}

Un processo è una macchina a stati. Ha una memoria interna, una ``casella di posta'' in entrata e una in uscita, e ad ogni passo legge la posta, aggiorna la memoria e magari spedisce messaggi.

\end{inbreve}

Formalmente un processo è un \textbf{automa} (I/O automaton, anche con infiniti stati):

\begin{itemize}
\tightlist
\item
  \(Q\): insieme degli stati interni; \(Q_{in} \subseteq Q\): stati iniziali possibili.
\item
  \(M\): insieme dei messaggi, nella forma \(\langle \text{mittente}, \text{destinatario}, \text{payload} \rangle\).
\item
  \(\text{InBuf}_j\): messaggi \textbf{arrivati} a \(p_j\) e non ancora letti. Vengono consumati al prossimo passo locale.
\item
  \(\text{OutBuf}_j\): messaggi \textbf{spediti} da \(p_j\) ma non ancora consegnati (``in volo'' sulla rete).
\end{itemize}

I buffer sono \textbf{multinsiemi}: lo stesso messaggio può comparire più volte e non c'è un ordine.

Lo \textbf{stato} completo di \(p_j\) è la terna \((q_j, \text{InBuf}_j, \text{OutBuf}_j)\).

\section{Esecuzioni asincrone}

\subsection{Lo scheduler (avversario)}

\begin{inbreve}

Chi decide ``cosa succede adesso'' non è nessun processo, ma un avversario immaginario. Gli algoritmi devono funzionare \textbf{qualunque} cosa scelga.

\end{inbreve}

Il sistema avanza un evento alla volta, e ogni volta un \textbf{avversario} (lo \emph{scheduler}) sceglie quale evento far succedere. Gli eventi possibili sono due:

\begin{itemize}
\tightlist
\item
  \(\text{Del}(i, j, m)\) (\emph{delivery}): la rete consegna \(m\), cioè lo sposta da \(\text{OutBuf}_i\) a \(\text{InBuf}_j\).
\item
  \(\text{Exec}(i)\): \(p_i\) fa \textbf{un passo} del suo automa: legge l'InBuf, cambia stato, eventualmente mette nuovi messaggi nell'OutBuf.
\end{itemize}

Pensare a uno scheduler cattivo serve a modellare il fatto che nella realtà non controlliamo i tempi: la rete può essere lenta, un processo può essere rallentato, ecc. Se l'algoritmo resiste al peggior scheduler, resiste a tutto.

\subsection{Configurazione, eventi abilitati, esecuzione}

\begin{itemize}
\tightlist
\item
  \textbf{Configurazione} \(C_t\): la ``fotografia'' del sistema al tempo \(t\), cioè il vettore degli stati di tutti i processi: \(C_t[j] = (q_j, \text{InBuf}_j, \text{OutBuf}_j)\).
\item
  Un evento è \textbf{abilitato} in \(C\) se può accadere in quella configurazione. \(\text{Exec}(i)\) è sempre abilitato; \(\text{Del}(i,j,m)\) è abilitato solo se \(m\) si trova davvero in \(\text{OutBuf}_i\) (non puoi consegnare un messaggio mai spedito).
\item
  \textbf{Esecuzione}: sequenza \textbf{infinita} che alterna configurazioni ed eventi
\end{itemize}

\[
E = (C_0, e_0, C_1, e_1, C_2, e_2, \dots)
\]

dove ogni \(e_t\) è abilitato in \(C_t\) e \(C_{t+1}\) è il risultato di applicare \(e_t\) a \(C_t\).

\textbf{Esempio}: \(C_0 \xrightarrow{\text{Exec}(1)} C_1 \xrightarrow{\text{Del}(1,2,m')} C_2 \xrightarrow{\text{Exec}(2)} C_3 \dots\) significa: \(p_1\) fa un passo e spedisce \(m'\); la rete consegna \(m'\) a \(p_2\); \(p_2\) fa un passo e lo legge.

\textbf{Asincrono} vuol dire che non c'è nessun limite di tempo: tra un evento e l'altro dello stesso processo, o tra l'invio e la consegna di un messaggio, possono passare quanti eventi si vuole.

\subsection{Esecuzione fair}

\begin{inbreve}

Un'esecuzione è fair (equa) se lo scheduler non ``dimentica'' nessuno per sempre: ogni processo prima o poi avanza, e ogni messaggio prima o poi arriva.

\end{inbreve}

Un'esecuzione \(E\) è \textbf{fair} se:

\begin{enumerate}
\def\labelenumi{\arabic{enumi}.}
\tightlist
\item
  ogni processo esegue \textbf{infiniti} passi locali, e
\item
  ogni messaggio spedito viene \textbf{prima o poi} consegnato da un evento \(\text{Del}\).
\end{enumerate}

\textbf{Perché serve?} Senza fairness lo scheduler potrebbe scegliere sempre e solo \(\text{Exec}(1)\): \(p_2\) non farebbe mai niente e i messaggi non arriverebbero mai. Nessun algoritmo potrebbe garantire che qualcosa ``prima o poi succede''. La fairness è l'ipotesi minima per poter parlare di progresso.

\textbf{Fairness con i guasti.} La definizione sopra è pensata per un sistema \textbf{senza guasti}. Con i crash non funziona: un processo crashato non viene più schedulato, quindi non fa infiniti passi e l'esecuzione non sarebbe mai fair. Per questo, con i guasti, la definizione si restringe ai processi corretti. Un'esecuzione è fair se:

\begin{enumerate}
\def\labelenumi{\arabic{enumi}.}
\tightlist
\item
  ogni processo \textbf{corretto} esegue infiniti passi locali, e
\item
  ogni messaggio spedito \textbf{da} un processo corretto \textbf{a} un processo corretto viene prima o poi consegnato.
\end{enumerate}

\section{Esecuzione locale e indistinguibilità}

\subsection{Esecuzione locale (local view)}

\begin{inbreve}

È quello che un singolo processo ``vede'' dell'esecuzione: solo gli eventi che lo riguardano, nell'ordine in cui sono avvenuti.

\end{inbreve}

\(E|p_j\) è la sottosequenza degli eventi di \(E\) che toccano \(p_j\) (i suoi passi e le consegne di messaggi a lui), nell'ordine originale.

\subsection{Indistinguibilità}

\begin{inbreve}

Due esecuzioni diverse possono sembrare identiche a ogni processo. Se succede, nessun processo può capire in quale delle due si trova.

\end{inbreve}

Esempio con 4 processi:

\[
E = (C_0,\ \text{Exec}(0),\ C_1,\ \text{Del}(0,1,m),\ C_2,\ \text{Exec}(3),\ C_3,\ \text{Del}(3,2,m'),\ \dots)
\]

\[
E' = (C_0,\ \text{Exec}(3),\ C'_1,\ \text{Del}(3,2,m'),\ C'_2,\ \text{Exec}(0),\ C'_3,\ \text{Del}(0,1,m),\ \dots)
\]

In \(E\) prima comunicano \(p_0 \to p_1\), poi \(p_3 \to p_2\); in \(E'\) al contrario. Ma le due coppie non interagiscono, quindi ogni processo vede esattamente la stessa storia: \(E|p_j = E'|p_j\) per ogni \(p_j\). Diciamo che \(E\) ed \(E'\) sono \textbf{indistinguibili}.

\textbf{Teorema}: in un sistema asincrono è impossibile ricostruire l'esecuzione esatta.

In parole: nessun processo sa davvero ``in che ordine sono successe le cose'' nel sistema globale. Questa idea (un processo non distingue due mondi diversi) è lo strumento principale per dimostrare che certi problemi sono \textbf{impossibili}.

\section{Sincrono vs asincrono}

Nel corso vedremo tre modelli di tempo:

\begin{itemize}
\tightlist
\item
  \textbf{Asincrono}: nessun limite sui tempi (quello visto sopra).
\item
  \textbf{Eventualmente sincrono}: asincrono per un po', poi da un certo momento (ignoto) diventa sincrono. Si vedrà nel modulo M2.
\item
  \textbf{Sincrono}: esistono limiti \textbf{fissi e noti} su (i) quanto ci mette un messaggio ad arrivare, (ii) quanto dura un'azione di un processo, (iii) quanto tempo passa tra due azioni.
\end{itemize}

Nel sincrono si può dire ``se non ricevo risposta entro 5 secondi, l'altro è morto''. Nell'asincrono no: un processo lento e uno guasto sono indistinguibili.

\section{Guasti dei processi}

\subsection{Crash-stop e bizantino}

\begin{inbreve}

Un processo si può rompere in due modi. \textbf{Crash}: si spegne e non fa più nulla. \textbf{Bizantino}: impazzisce o diventa malevolo e fa qualunque cosa, anche mentire.

\end{inbreve}

\begin{itemize}
\tightlist
\item
  \(\text{Crash}(p_j)\): da quel momento \(p_j\) non esegue più passi (niente più \(\text{Exec}(j)\)). Semplicemente si ferma.
\item
  \(\text{Byz}(p_j)\): da quel momento \(p_j\) si comporta in modo \textbf{arbitrario}: quando fa un passo non segue più il suo automa. Può mandare messaggi falsi, messaggi diversi a processi diversi, tacere, ecc. È il modello giusto per un attaccante che ha preso il controllo di una macchina.
\end{itemize}

I guasti bizantini \textbf{contengono} i crash: fermarsi è uno dei possibili comportamenti arbitrari. Quindi un algoritmo che tollera i bizantini tollera anche i crash, non viceversa.

\subsection{Processi corretti e il parametro f}

\begin{itemize}
\tightlist
\item
  Un processo è \textbf{corretto} se non subisce mai guasti. Attenzione: durante l'esecuzione \textbf{nessuno sa} chi è corretto.
\item
  \(f\) = numero massimo di processi che possono guastarsi in un'esecuzione. Di solito lo si limita, ad esempio \(f \le n-1\) (almeno uno sopravvive).
\item
  Non si fanno \textbf{mai} ipotesi su \textbf{quando} avviene un guasto: può capitare in qualunque momento, anche a metà algoritmo.
\end{itemize}

\section{Astrazioni e specifiche}

\subsection{Come si costruisce un'astrazione distribuita}

\begin{inbreve}

Un'astrazione è un ``oggetto'' con un comportamento promesso (es. ``un canale che non perde messaggi''). Prima si scrive cosa promette, poi si costruisce, poi si dimostra che mantiene la promessa.

\end{inbreve}

Il metodo del corso in 4 passi:

\begin{enumerate}
\def\labelenumi{\arabic{enumi}.}
\tightlist
\item
  Definire il \textbf{modello} del sistema (processi, tempi, guasti).
\item
  Definire il problema come \textbf{astrazione}: un'interfaccia più un elenco di \textbf{proprietà} scritte senza ambiguità.
\item
  \textbf{Implementare} l'astrazione con un protocollo distribuito.
\item
  \textbf{Dimostrare} che il protocollo soddisfa tutte le proprietà.
\end{enumerate}

Ogni astrazione comunica con chi la usa tramite \textbf{eventi}:

\begin{itemize}
\tightlist
\item
  \textbf{Request}: chiamate dall'alto verso l'astrazione (es. \(\langle \text{Send} \mid q, m \rangle\): ``spedisci \(m\) a \(q\)'').
\item
  \textbf{Indication}: notifiche dall'astrazione verso l'alto (es. \(\langle \text{Deliver} \mid p, m \rangle\): ``ti è arrivato \(m\) da \(p\)'').
\end{itemize}

\subsection{Safety e liveness}

\begin{inbreve}

Ogni proprietà è di uno di due tipi. \textbf{Safety}: ``non succede mai niente di male''. \textbf{Liveness}: ``prima o poi succede qualcosa di buono''.

\end{inbreve}

\textbf{Safety}

\begin{itemize}
\tightlist
\item
  Informale: \emph{qualcosa di cattivo non succede mai}.
\item
  Se viene violata al tempo \(t\), resta violata per sempre: non si può più rimediare.
\item
  Formale: se \(E\) viola la proprietà, esiste un \textbf{prefisso finito} \(E'\) di \(E\) tale che \textbf{ogni} estensione di \(E'\) la viola.
\end{itemize}

\textbf{Liveness}

\begin{itemize}
\tightlist
\item
  Informale: \emph{qualcosa di buono prima o poi succede}.
\item
  Non può essere violata in tempo finito: c'è sempre speranza.
\item
  Formale: ogni esecuzione \textbf{finita} che non la soddisfa si può \textbf{estendere} in una che la soddisfa.
\end{itemize}

\textbf{Trucco pratico} per classificarle: chiediti ``posso indicare il momento esatto in cui la proprietà è stata rotta?''. Se sì, è safety. Se la violazione si vede solo ``all'infinito'' (aspetti per sempre e non succede), è liveness.

\textbf{Esempi}

\begin{itemize}
\tightlist
\item
  \emph{Mutua esclusione}: se \(p\) ha la risorsa \(r\) al tempo \(t\), nessun altro ce l'ha al tempo \(t\). → \textbf{safety} (il momento in cui due la tengono insieme è la violazione).
\item
  \emph{No-deadlock}: se \(r\) è libera, prima o poi qualcuno la ottiene. → \textbf{liveness}.
\item
  \emph{No-starvation}: se \(p\) chiede \(r\), prima o poi la ottiene. → \textbf{liveness}.
\item
  \emph{``Al massimo 7 consegne''}: ogni messaggio viene consegnato al più 7 volte. → \textbf{safety}: all'ottava consegna è violata e nessuna estensione rimedia.
\end{itemize}

\subsection{Proprietà scritte male}

Una buona specifica separa safety e liveness in proprietà distinte. Esempio scritto male:

\begin{nota}

``Se \(p\) manda \(m\) a \(q\), allora \(q\) prima o poi lo consegna e la consegna è unica.''

\end{nota}

Mescola due cose. Meglio:

\begin{itemize}
\tightlist
\item
  (liveness) Se \(p\) manda \(m\) a \(q\), \(q\) prima o poi consegna \(m\).
\item
  (safety) Se \(p\) manda \(m\) a \(q\), \(m\) è consegnato al più una volta.
\end{itemize}

\section{I link: dalla rete inaffidabile al canale perfetto}

L'obiettivo: partire da una rete realistica (che perde e duplica messaggi) e costruire, strato sopra strato, canali sempre più affidabili.

\begin{verbatim}
FIFO perfect link      ← consegna tutto, una volta, in ordine
       ▲
Perfect link (pl)       ← consegna tutto, una volta
       ▲
Stubborn link (sl)      ← consegna tutto, infinite volte
       ▲
Fair-loss link (fll)    ← la rete reale
\end{verbatim}

\subsection{Fair-loss link}

\begin{inbreve}

È il modello della rete vera. Può perdere messaggi, ma se insisti a rimandare lo stesso messaggio, prima o poi passa. Può duplicare, ma solo un numero finito di volte. Non inventa messaggi.

\end{inbreve}

\textbf{Prima di tutto: cos'è un link come astrazione.} Pensalo come una scatola nera tra due processi, con un ingresso e un'uscita:

\begin{verbatim}
  p                                             q
  │  ⟨Send | q, m⟩                ⟨Deliver | p, m⟩  ▲
  ▼                                                 │
┌───────────────────── LINK ──────────────────────┐
│  qui dentro la rete può perdere o copiare m      │
└─────────────────────────────────────────────────┘
\end{verbatim}

\begin{itemize}
\tightlist
\item
  \textbf{Send} è la richiesta del mittente: ``prova a far arrivare \(m\) a \(q\)''.
\item
  \textbf{Deliver} è l'evento lato destinatario: ``il link ti consegna \(m\), arrivato da \(p\)''.
\end{itemize}

Un Send \textbf{non garantisce} un Deliver. Tutte le proprietà del link sono regole su che rapporto c'è tra i Send fatti da una parte e i Deliver che escono dall'altra.

\textbf{Perché proprio questo modello.} È la rete reale vista dal basso (tipo UDP): i pacchetti si perdono (congestione, buffer pieni), a volte arrivano doppi (un router o un livello sotto ritrasmette), ma la rete non si inventa pacchetti dal nulla. L'analogia è un \textbf{postino distratto}: certe lettere le perde, certe le consegna due volte, ma non scrive lettere di suo pugno. E se continui a spedire la stessa lettera, prima o poi qualcuna arriva.

Proprietà, una per una:

\textbf{FL1 (Fair-loss)}: se un processo corretto \(p\) manda \(m\) \textbf{infinite volte} a un processo corretto \(q\), allora \(q\) consegna \(m\) infinite volte.

\begin{itemize}
\tightlist
\item
  Dice cosa \textbf{non} può fare la rete: perdere \emph{per sempre} un messaggio che continui a rispedire. Può perdere molte copie, ma non tutte.
\item
  La formulazione con ``infinite volte'' sembra strana, ma è l'unica onesta: la rete non promette nulla su \textbf{un singolo} invio (ogni copia può andare persa). Se ogni copia si perde con probabilità \(pr < 1\), insistendo all'infinito ne passano infinite.
\item
  Perché ``corretti''? Se \(p\) crasha smette di rispedire; se \(q\) crasha non consegna più nulla. La promessa vale solo se entrambi restano vivi.
\item
  Tipo: \textbf{liveness} (promette che qualcosa di buono, le consegne, continuerà ad accadere).
\end{itemize}

\textbf{FL2 (Finite duplication)}: se \(p\) manda \(m\) un numero \textbf{finito} di volte, \(q\) non può consegnarlo infinite volte.

\begin{itemize}
\tightlist
\item
  La rete può creare copie, ma un numero limitato: se mandi \(m\) 3 volte, magari ne arrivano 5, ma non infinite.
\item
  Senza FL2 la rete potrebbe ``spammare'' \(q\) all'infinito con una sola spedizione, e soprattutto potrebbe far sembrare che \(p\) stia ritrasmettendo quando invece si è fermato.
\item
  Tipo: \textbf{liveness}, ed è il caso furbo. Dato un pezzo finito di esecuzione con, diciamo, 1000 consegne di \(m\), non hai violato niente: 1000 è finito, e l'esecuzione si può sempre continuare senza altre consegne. La violazione (``infinite consegne'') si vede solo all'infinito, quindi per definizione non è safety.
\end{itemize}

\textbf{FL3 (No creation)}: se \(q\) consegna \(m\) con mittente \(p\), allora \(p\) ha davvero spedito \(m\) a \(q\).

\begin{itemize}
\tightlist
\item
  La rete non inventa messaggi, non li modifica e non cambia il mittente.
\item
  Tipo: \textbf{safety}. Se \(q\) consegna un messaggio mai spedito, c'è un istante preciso in cui è successo, e da lì la proprietà è rotta per sempre.
\end{itemize}

\textbf{Esempio.} \(p\) manda \(m\) tre volte e poi si ferma.

\begin{itemize}
\tightlist
\item
  Consentito: nessuna consegna (tutte perse); 1 consegna; 4 consegne (una copia duplicata).
\item
  Vietato: infinite consegne (viola FL2); una consegna di \(m'\) che \(p\) non ha mai mandato (viola FL3).
\end{itemize}

Se invece \(p\) continua a mandare \(m\) per sempre, ``zero consegne'' diventa vietato (viola FL1).

\textbf{Cosa il fair-loss NON garantisce}: che un singolo invio arrivi, che non ci siano duplicati, l'ordine di arrivo, i tempi. Il resto della lezione serve a costruire queste garanzie sopra il fair-loss: lo \textbf{stubborn} sfrutta FL1 rispedendo per sempre, il \textbf{perfect} elimina i duplicati, il \textbf{FIFO} sistema l'ordine.

\subsection{Come si scrivono gli algoritmi}

Lo pseudocodice del corso è \textbf{a eventi}: un algoritmo è un insieme di \textbf{handler} del tipo \texttt{upon\ event\ ⟨...⟩\ do\ ...}. Ogni handler reagisce a un evento (una richiesta dall'alto, una consegna dal basso, un timer) e può generare nuovi eventi con \texttt{trigger}.

Gli handler sono \textbf{atomici}: un handler, una volta iniziato, finisce senza essere interrotto (tranne da un crash). Questo è importante nelle dimostrazioni.

\subsection{Stubborn link}

\begin{inbreve}

Un canale ``testardo'': tutto quello che mandi arriva, anzi arriva infinite volte. Si ottiene rimandando ogni messaggio per sempre.

\end{inbreve}

Proprietà:

\begin{itemize}
\tightlist
\item
  \textbf{SL1 (Stubborn delivery)}: se un processo corretto \(p\) manda \(m\) a un processo corretto \(q\), \(q\) consegna \(m\) infinite volte.
\item
  \textbf{FL3 (No creation)}: come sopra.
\end{itemize}

Algoritmo \emph{Retransmit Forever} (usa un fair-loss link \texttt{fll}):

\begin{verbatim}
upon event ⟨sl, Init⟩ do
    sent := ∅
    starttimer(Δ)

upon event ⟨Timeout⟩ do            // ogni Δ: rimanda tutto
    forall (q, m) ∈ sent do
        trigger ⟨fll, Send | q, m⟩
    starttimer(Δ)

upon event ⟨sl, Send | q, m⟩ do
    trigger ⟨fll, Send | q, m⟩
    sent := sent ∪ {(q, m)}

upon event ⟨fll, Deliver | p, m⟩ do
    trigger ⟨sl, Deliver | p, m⟩
\end{verbatim}

\textbf{Perché funziona}

\begin{itemize}
\tightlist
\item
  \emph{No creation}: lo stubborn consegna un messaggio solo quando il fair-loss sotto lo consegna. Se consegnasse un messaggio mai spedito, l'avrebbe inventato il fair-loss, che però per FL3 non lo fa.
\item
  \emph{Stubborn delivery}: \(m\) resta in \texttt{sent} per sempre, quindi il timer lo rimanda \textbf{infinite volte}. Per FL1 il fair-loss lo consegna infinite volte, e ogni consegna del fair-loss diventa una consegna dello stubborn.
\end{itemize}

\subsection{Perfect link}

\begin{inbreve}

Il canale che vorremmo sempre: ogni messaggio arriva, arriva \textbf{una volta sola}, e non ne arrivano di inventati. Si ottiene dallo stubborn buttando via i duplicati.

\end{inbreve}

Proprietà:

\begin{itemize}
\tightlist
\item
  \textbf{PL1 (Reliable delivery)}: se \(p\) corretto manda \(m\) a \(q\) corretto, \(q\) prima o poi consegna \(m\). → liveness
\item
  \textbf{PL2 (No duplication)}: ogni messaggio è consegnato al più una volta. → safety
\item
  \textbf{FL3 (No creation)}. → safety
\end{itemize}

Algoritmo \emph{Eliminate Duplicates} (usa uno stubborn link \texttt{sl}):

\begin{verbatim}
upon event ⟨pl, Init⟩ do
    delivered := ∅

upon event ⟨pl, Send | q, m⟩ do
    trigger ⟨sl, Send | q, m⟩

upon event ⟨sl, Deliver | p, m⟩ do
    if m ∉ delivered then           // filtro anti-duplicati
        delivered := delivered ∪ {m}
        trigger ⟨pl, Deliver | p, m⟩
\end{verbatim}

\textbf{Perché funziona}

\begin{itemize}
\tightlist
\item
  \emph{No creation}: come prima, si eredita dallo stubborn.
\item
  \emph{No duplication}: supponiamo \(m\) consegnato due volte, ai tempi \(t < t'\). Al tempo \(t\) l'handler (atomico) aggiunge \(m\) a \texttt{delivered}. Al tempo \(t'\) quindi \(m \in\) \texttt{delivered} e l'\texttt{if} blocca la consegna. Contraddizione.
\item
  \emph{Reliable delivery}: lo stubborn consegna \(m\) (anzi infinite volte). La prima volta \(m\) non è in \texttt{delivered}, quindi viene consegnato.
\end{itemize}

\textbf{Attenzione (correzione nelle slide)}: una prima versione di PL1 diceva solo ``se \(p\) corretto manda \(m\) a \(q\), \(q\) consegna \(m\)''. Con i crash è sbagliata: se \(q\) crasha non può consegnare niente. Serve che \textbf{anche} \(q\) \textbf{sia corretto}. In generale, nelle proprietà di liveness bisogna sempre dire chi deve essere corretto.

\subsection{FIFO perfect link}

\begin{inbreve}

Un perfect link che in più rispetta l'ordine: i messaggi arrivano nello stesso ordine in cui sono stati spediti.

\end{inbreve}

Proprietà: PL1, PL2, FL3 e in più

\begin{itemize}
\tightlist
\item
  \textbf{FIFO}: se \(q\) consegna \(m\) prima di \(m'\) (entrambi da \(p\)), allora \(p\) ha spedito \(m\) prima di \(m'\).
\end{itemize}

\textbf{Idea}: numerare i messaggi. Il mittente attacca a ogni messaggio un numero di sequenza crescente; il destinatario tiene da parte i messaggi arrivati ``in anticipo'' e li consegna solo quando arriva quello col numero atteso.

\begin{verbatim}
upon event ⟨fpl, Init⟩ do
    pending := ∅
    forall q ∈ Π do
        seq[q]  := 0      // ultimo numero usato verso q
        next[q] := 1      // prossimo numero atteso da q

upon event ⟨fpl, Send | q, m⟩ do
    seq[q] := seq[q] + 1
    trigger ⟨pl, Send | q, [seq[q], m]⟩

upon event ⟨pl, Deliver | s, [k, m]⟩ do
    pending := pending ∪ {(s, k, m)}

upon exists (s, k, m) ∈ pending such that k = next[s] do
    pending := pending \ {(s, k, m)}
    next[s] := next[s] + 1
    trigger ⟨fpl, Deliver | s, m⟩
\end{verbatim}

\textbf{Perché funziona}: PL1, PL2, FL3 vengono direttamente dal perfect link sotto. Per FIFO: \texttt{next{[}s{]}} cresce di 1 a ogni consegna, quindi i messaggi di \(s\) vengono consegnati esattamente nell'ordine 1, 2, 3, \ldots{} dei loro numeri, che è l'ordine di invio. Se \(m'\) (numero più alto) venisse consegnato prima di \(m\), \texttt{next} avrebbe dovuto già passare per il numero di \(m\), cioè \(m\) sarebbe già stato consegnato. Contraddizione.

Nota: i numeri di sequenza risolvono anche un altro problema, cioè due messaggi con lo \textbf{stesso contenuto} (vedi l'esercizio del contatore sotto).

\section{Esercizi e homework}

\begin{approfondimento}{Homework: safety o liveness? (soluzione proposta, non nelle slide)}

\begin{itemize}
\tightlist
\item
  \textbf{Fairness di un'esecuzione} → liveness: qualunque prefisso finito si può estendere in un'esecuzione fair, non si può mai dire ``ormai è rotta''.
\item
  \textbf{Al più f guasti} → safety: al momento del guasto numero \(f+1\) la proprietà è violata per sempre.
\item
  \textbf{Eventual recovery} (chi crasha prima o poi riparte) → liveness.
\item
  \textbf{Time-bounded recovery} (se c'è un recovery, avviene entro \(k\) eventi dal guasto) → safety: un recovery arrivato dopo più di \(k\) eventi è una violazione visibile in un prefisso finito.
\item
  \textbf{Time-bounded implica eventual?} No.~Così come è scritta, se il processo crasha e \textbf{non} riparte mai, time-bounded è soddisfatta ``a vuoto'' (non c'è nessun recovery da controllare), mentre eventual è violata.
\item
  \textbf{Riscrivere la proprietà sugli insiemi}: separare in (liveness) ``se un processo corretto produce un output, prima o poi ogni processo corretto produce un output'' e (safety) ``per due output qualunque \(O_z, O_t\) vale \(O_z \subseteq O_t\) oppure \(O_t \subseteq O_z\)''.
\end{itemize}

\end{approfondimento}

\begin{approfondimento}{Homework 1: stubborn senza `corretto' (soluzione nelle slide)}

Se SL1 dice solo ``se \(p\) manda \(m\) a \(q\), \(q\) lo consegna infinite volte'', senza richiedere che \(p\) sia corretto, l'algoritmo non funziona. Se \(p\) crasha smette di ritrasmettere, quindi manda \(m\) sul fair-loss solo un numero finito di volte, e il fair-loss può perderle tutte. Lo stesso argomento mostra che \textbf{nessun} algoritmo sopra un fair-loss può soddisfare quella proprietà.

\textbf{Domanda (*)}: con memoria limitata a \(k\) messaggi si può fare un perfect link sopra un fair-loss? \emph{Idea (proposta, da verificare a lezione)}: sembra di no. Per non consegnare duplicati il destinatario deve ricordare quali messaggi ha già consegnato, e una copia vecchia può arrivare in ritardo arbitrario. Con memoria limitata prima o poi deve ``dimenticare'' qualcosa e non riesce più a riconoscere un duplicato.

\end{approfondimento}

\begin{approfondimento}{Homework 2: l'algoritmo di Caio senza timer (soluzione nelle slide)}

Caio propone uno stubborn senza timer: a ogni nuova richiesta di invio rimanda \textbf{tutti} i messaggi accumulati in \texttt{set}.

\begin{verbatim}
upon event ⟨sl, Send | m, q⟩ do
    set := set ∪ {(m, q)}
    forall (x, y) ∈ set do
        trigger ⟨fll, Send | x, y⟩

upon event ⟨fll, Deliver | m, q⟩ do
    trigger ⟨sl, Deliver | m, q⟩
\end{verbatim}

\textbf{Cosa non va}: i messaggi vengono rimandati solo quando arriva una nuova richiesta di invio. Se \(p\) manda \(m\) e poi non manda più niente, \(m\) va sul fair-loss una volta sola e può andare perso.

\textbf{Quando funziona}: solo nelle esecuzioni in cui \(p\) continua a mandare infiniti messaggi.

\textbf{Perché è impossibile senza timer}: per garantire la consegna bisogna mandare \(m\) sul fair-loss infinite volte. Senza timer (o ciclo infinito, che è la stessa cosa) e senza infiniti input dall'esterno, dopo un certo numero di passi \(p\) non cambia più stato e non fa più nulla. Quindi esiste un'esecuzione in cui \(m\) viene mandato un numero finito di volte, e il fair-loss può perderle tutte.

\end{approfondimento}

\begin{approfondimento}{Esercizio: il server contatore (soluzione nelle slide + commento mio)}

\textbf{Testo}: un server ha una variabile \(v = 0\) e accetta INC, DEC e GET (risponde con \(v\)). Comunica col client tramite il perfect link visto sopra. Il client manda INC, INC, GET e riceve \textbf{1} invece di 2. Rimanda GET e non riceve più risposta. Perché?

\textbf{Spiegazione 1: ordine.} Il perfect link non garantisce l'ordine: il server può ricevere INC, GET, INC e rispondere 1.

\textbf{Spiegazione 2: duplicati.} Anche se l'ordine è giusto: il filtro \texttt{if\ m\ ∉\ delivered} guarda il \textbf{contenuto} del messaggio. Il secondo INC è identico al primo, quindi viene scartato come duplicato, e \(v\) vale 1. Il secondo GET è identico al primo, quindi viene scartato anche lui: ecco perché non arriva più nessuna risposta. (L'indizio ``quanto aspetteresti?'' punta a questo: nell'asincrono non puoi distinguere un messaggio lento da uno buttato via.)

\textbf{Soluzione}: dare a ogni messaggio un \textbf{identificatore unico} (es. numero di sequenza), così due INC sono messaggi diversi; per l'ordine usare il FIFO link, che fa entrambe le cose.

\end{approfondimento}

\section{Riepilogo}

\begin{itemize}
\tightlist
\item
  \textbf{Modello}: \(n\) processi = automi con InBuf/OutBuf, collegati da link; uno \textbf{scheduler avversario} sceglie gli eventi \(\text{Exec}(i)\) e \(\text{Del}(i,j,m)\).
\item
  \textbf{Esecuzione} = sequenza infinita configurazione/evento; \textbf{fair} = ogni processo avanza infinite volte e ogni messaggio arriva.
\item
  \textbf{Indistinguibilità}: esecuzioni diverse possono sembrare uguali a tutti i processi → in asincrono non si ricostruisce l'esecuzione esatta.
\item
  \textbf{Tempo}: asincrono / eventualmente sincrono / sincrono.
\item
  \textbf{Guasti}: crash-stop (si ferma) ⊂ bizantino (fa di tutto); al più \(f\) guasti, in momenti imprevedibili.
\item
  \textbf{Safety} = niente di male, violazione visibile in tempo finito; \textbf{liveness} = prima o poi qualcosa di buono, mai violabile in tempo finito.
\item
  \textbf{Link}: fair-loss → (ritrasmetti per sempre) → stubborn → (scarta duplicati) → perfect → (numeri di sequenza) → FIFO perfect.
\end{itemize}
